\section{Related Work}
\label{sec:related}
\textbf{Fashion recognition and retrieval.} DeepFashion~\cite{liu2016deepfashion} and
DeepFashion2~\cite{ge2019deepfashion2} provide large benchmarks for clothing categories,
attributes, landmarks and consumer-to-shop retrieval, and street-to-shop matching finds the exact
product shown in a photo~\cite{kiapour2015where}. Clothing style has been rated by crowds and
classified~\cite{kiapour2014hipster} and learned as embeddings from weak
labels~\cite{simoserra2016fashion}. These systems recognize \emph{what} a garment
is or find the same product online; they do not predict a second-hand item's provenance (house and
production era), which is what drives the price of archive pieces.
\textbf{Dating images and fashion.} Estimating when a photo was taken has been studied for
historical colour photographs~\cite{palermo2012dating} and a century of yearbook
portraits~\cite{ginosar2017century}. Closest to our era task,
\citet{vittayakorn2017when} estimate when clothing was made from dated Flickr and museum images, and
Runway to Realway~\cite{vittayakorn2015runway} estimates the brand, season and year of runway looks;
trend studies model how styles evolve from product and outfit
photos~\cite{alhalah2017fashion,he2016ups,mall2019geostyle,hsiao2021culture}. These works use the
whole photo, garment or outfit; we target label-level cues for specific houses, where printed date
codes (and, for Margiela, label changes) give sharp era boundaries.
\textbf{Price prediction.} \citet{han2020price} regress the prices of second-hand marketplace items
from images and text. We instead estimate value by retrieval, which is how resellers price archive
pieces (by finding comparable sales) and which returns the comparables themselves.
\textbf{Representations and robustness.} We build on pretrained encoders: CLIP~\cite{radford2021clip},
DINOv2~\cite{oquab2024dinov2}, and ImageNet~\cite{deng2009imagenet} ResNets~\cite{he2016resnet};
we follow standard transfer protocols~\cite{kornblith2019transfer}, and our schedule loosely follows
the observation that training a linear head before fine-tuning preserves robustness under
distribution shift~\cite{kumar2022finetuning}. Telling brands apart is fine-grained
recognition~\cite{wei2022finegrained}. We evaluate with common
corruptions~\cite{hendrycks2019corruptions} and check for shortcut
learning~\cite{geirhos2020shortcut}.
\textbf{Reading labels.} Brand marks are close to logo detection~\cite{wang2022logodet3k}, and
printed label text calls for OCR: scene-text detection and recognition~\cite{baek2019craft,shi2017crnn}
and transformer recognizers~\cite{li2023trocr}.
